%%=============================================================================
%% Methodologie
%%=============================================================================
\chapter{\IfLanguageName{dutch}{Methodologie}{Methodology}}%
\label{ch:methodologie}

Deze bachelorproef werd uitgewerkt in vijf opeenvolgende fases. Elke fase had een duidelijke doelstelling met eigen deliverables en een specifieke techniek. Zo werd de focus telkens op een ander deel van het probleem gelegd en werd er systematisch verder gebouwd aan een conclusie. In totaal werden er 14 werkdagen voorzien, waarvan 1 werkdag per werkweek besteed werd aan het onderzoek. Een overzicht van de tijdsbesteding per fase wordt weergegeven in Figuur~\ref{fig:gantt}.

\section{Fase 1: Analyse van het probleemdomein}

In de eerste fase van dit onderzoek werd het probleemdomein verder onderzocht, door middel van het huidige evaluatieproces binnen de opleidingsonderdelen Front-End Web Development en Web Services te analyseren. Hiervoor werden er interviews afgenomen met alle leden van de doelgroep, zijnde de 3 lectoren die beide vakken geven.

\vspace{1em}

Het doel van deze fase was:
\begin{itemize}
    \item verder verdiepen binnen het probleem, de evaluatiecriteria en de methoden die de lectoren momenteel hanteren.
    \item een gemiddelde tijdsbesteding per project in kaart brengen
    \item ervaringen met inconsistenties binnen beoordelingen te weten komen
\end{itemize}

Het resultaat van deze fase bestond uit een zeer duidelijke nood aan onderzoek, een concreet inzicht in de richting die het onderzoek uit zou moeten gaan, een overzicht van de verschillende evaluatiecriteria en hoe de evaluatie van studentenprojecten verloopt, en een verzameling van projecten met uiteenlopende beoordelingen, om te kunnen gebruiken als testprojecten voor bij een volgende fase.

\section{Fase 2: Literatuurstudie}

De tweede fase bestond uit een systematische literatuurstudie rond:
\begin{itemize}
    \item autograding binnen programmeeronderwijs
    \item het gebruik van Large Language Models (LLM's) voor code-analyse
    \item betrouwbaarheid, hallucinaties en bias in LLM-output
    \item prompt-engineering technieken
\end{itemize}

Uit deze literatuurstudie werden verschillende relevante inzichten rondom het inzetten van LLM's verworven die voor de technische uitwerking van pas komen. Op basis van de technische vereisten werd gekozen voor de inzet van \textbf{qwen3-coder:480B Cloud}. Dit model werd geselecteerd vanwege zijn sterke redeneercapaciteiten, het feit dat het als cloud-model geen lokale rekenkracht vereist, en de gratis beschikbaarheid voor dit onderzoek.

\section{Fase 3: Ontwerp en opzetten testomgeving en PoC}

In de derde fase werd een reproduceerbare testomgeving ontwikkeld waarin het geselecteerde LLM systematisch ingezet kon worden voor de analyse van studentenprojecten. Dit omvat:
\begin{itemize}
    \item het implementeren van een opsplitsingsmechanisme om rekening te houden met contextlimieten
    \item het ontwikkelen van scripts voor automatische projectanalyse
    \item het ontwerpen van een gestandaardiseerde promptstructuur
    \item het definiëren van output standaarden
\end{itemize}

De PoC werd stapgewijs opgebouwd, uitgebreid en verbeterd, waardoor deze fase een grondig gedocumenteerde en reproduceerbare testomgeving oplevert als deliverable.

\section{Fase 4: Kwantitatieve evaluatie van de PoC}

In deze fase werd de werking van de PoC geëvalueerd op basis van:

\begin{itemize}
    \item correctheid van criteriadetectie door het LLM
    \item consistentie van de AI-beoordelingen met de evaluaties van lectoren
    \item de kwaliteit van de gegenereerde feedback
\end{itemize}

De resultaten van de automatische analyse werden getoetst aan de hand van de eerder verzamelde handmatige beoordelingen om de betrouwbaarheid van de tool vast te stellen.

% =================================================================

\section{Fase 5: Conclusie}

In de laatste fase werden de bevindingen uit alle voorgaande fasen samengebracht. 

Hier werd beoordeeld:
\begin{itemize}
    \item in welke mate LLM’s bruikbaar zijn als ondersteuning bij evaluatie van studentenprojecten
    \item op welke criteria LLM’s sterker of zwakker presteren
    \item welke randvoorwaarden noodzakelijk zijn voor verantwoord gebruik binnen het kader van de AI Act
\end{itemize}

